\subsection{Regularität und Tangentialvektor}

\begin{definition}{}
(a) Eine Kurve $\gamma : I \to \mathbb{R}^n$ heißt \emph{stetig differenzierbar}, wenn alle Komponentenfunktionen $\gamma_i : I \to \mathbb{R}$ stetig differenzierbar sind. Wir nennen


\[
\gamma'(t) = \bigl( \gamma_1'(t),\, \gamma_2'(t),\, \dots,\, \gamma_n'(t) \bigr)
\]


den Tangentialvektor der Kurve.

(b) $\gamma$ heißt \emph{regular}, wenn $\gamma'(t) \neq 0$ für alle $t \in I$.

(c) Die Funktion 


\[
v(t) = \|\gamma'(t)\|
\]


heißt die \emph{Geschwindigkeit} der Kurve.
\end{definition}

\begin{beispielbox}
\begin{enumerate}
    \item \textbf{Gerade:} Sei 
    

\[
    \gamma(t) = a + t\,v,
    \]


    wobei $v \neq 0$. Dann ist $\gamma$ regelmäßig.

    \item \textbf{Kreis:} Eine Kreis-Kurve in $\mathbb{R}^2$ mit Mittelpunkt $a$ und Radius $r>0$ wird durch
    

\[
    \gamma(t) = a + r\,(\cos t,\, \sin t)
    \]


    beschrieben. Es gilt
    

\[
    \gamma'(t) = r\,(-\sin t,\, \cos t)
    \]


    und somit ist $v(t)=r>0$, also ist die Kurve regelmäßig.

    \item \textbf{Helix:} Definiere 
    

\[
    \gamma(t) = \bigl( r\cos t,\, r\sin t,\, c\,t \bigr),
    \]


    wobei $r>0$ und $c\neq 0$. Dann folgt
    

\[
    \gamma'(t) = \bigl( -r\sin t,\, r\cos t,\, c \bigr).
    \]


    Damit ist 
    

\[
    v(t)=\sqrt{r^2+c^2}>0,
    \]


    sodass die Helix regelmäßig ist.

    \item \textbf{Graph einer Funktion:} Der Graph einer stetigen Funktion $f : I \to \mathbb{R}$, gegeben durch
    

\[
    \gamma(t) = (t, f(t)),
    \]


    ist genau dann regelmäßig, wenn $\gamma'(t) = (1, f'(t)) \neq (0,0)$.

    \item \textbf{Neilsche Parabel:} Betrachte die Kurve
    

\[
    \gamma(t) = (t^3, \, t^3).
    \]


    Es gilt
    

\[
    \gamma'(t) = (3t^2,\, 3t^2).
    \]


    Insbesondere ist $\gamma'(0)=(0,0)$, sodass die Kurve bei $t=0$ nicht regelmäßig ist.
\end{enumerate}
\end{beispielbox}

\begin{bemerkungbox}
In einem regulären Punkt einer Kurve $\gamma : I \to \mathbb{R}^n$, d.h. für ein $t_0\in I$ mit $\gamma'(t_0)\neq 0$, ist die Tangente an die Kurve durch


\[
\ell = \gamma(t_0) + \mathbb{R}\,\gamma'(t_0)
\]


eindeutig bestimmt, wobei gilt


\[
\gamma'(t_0) = \lim_{h\to 0} \frac{\gamma(t_0+h)-\gamma(t_0)}{h}\,.
\]


\end{bemerkungbox}
